\chapter{Introduction}
\section{Abstract}
EDADSL (Electronic Design Automation Domain Specific Language) is a Domain-specific Language (DSL) specialized for programmatic Electronic Design Automation (EDA) and Hardware-Synthesis. 
It is based on an Electrical System and a constraint-driven optimisation engine.
This Document is a full outline and as of now the authoritative Specification of the EDADSL Language as well as the systems underlying it.

\section{Reading this Document}
This document is organized into the following parts:

\begin{itemize}
\item \textbf{Part I: Core Language} -- Defines the programming language constructs including lexical structure, types, variables, expressions, constants, control flow, and functions.
\item \textbf{Part II: Formal Semantics} -- Describes the evaluation model, execution order, constraint resolution, and type system rules.
\item \textbf{Part III: Core Formal Model} -- Defines the structure of programs, the electrical system, and the solver.
\item \textbf{Part IV: Hardware Model} -- Describes the circuit model, electrical objects, and querying system.
\item \textbf{Part V: Physical System} -- Defines physical constraints and the solver model.
\item \textbf{Part VI: Toolchain} -- Describes system flags, the compiler pipeline, code generation, standard library, external library integration, and debugging.
\item \textbf{Part VII: Reference} -- Contains exceptions, error model, future improvements, implementations, and references.
\item \textbf{Appendices} -- Contains standard libraries, examples, and diagrams.
\end{itemize}

Readers new to EDADSL should start with Parts I and IV, which define the language and electrical system. Those interested in physical layout should read Part V. Implementers should consult Parts II and III for formal specifications.


\part{Core Language}
\chapter{Design Principles}
\section{Design Goals}
EDADSL was designed with the following primary goals:

\begin{enumerate}
\item \textbf{Programmatic PCB Design} -- Enable engineers to describe circuits and layouts programmatically, enabling version control, automation, and reproducibility.

\item \textbf{Constraint-Driven Layout} -- Allow physical constraints (clearance, creepage, alignment) to be specified declaratively, with the solver handling the complex optimization.

\item \textbf{Electrical Awareness} -- Maintain a rich electrical model (parts, pins, connections, nets) that understands circuit topology, not just geometry.

\item \textbf{KiCad Integration} -- Generate industry-standard KiCad files (schematics and PCBs) for further editing and manufacturing.

\item \textbf{Readability} -- Use clear, prefixed syntax (elec.*, phys.*, syst.*) to make code self-documenting and easy to understand.

\item \textbf{Comprehensiveness} -- Provide a comprehensive set of built-in functions covering trigonometry, linear algebra, signal processing, and other domains needed for EDA workflows.

\item \textbf{Accessibility} -- Serve both professional engineers and hobbyists with appropriate defaults while supporting advanced features for complex designs.
\end{enumerate}

\section{Applications and Application-Areas with special Priority during Design}
EDADSL is particularly well-suited for:

\begin{itemize}
\item \textbf{Power Electronics} -- Where creepage distances, high-current traces, and thermal management are critical. The constraint system supports isolation requirements between high-voltage and low-voltage sections.

\item \textbf{High-Speed Digital Design} -- Where trace length matching, impedance control, and signal integrity are important. The \texttt{phys.connmatch} constraint supports length and impedance matching.

\item \textbf{Mixed-Signal Designs} -- Where analog and digital sections must be carefully isolated. The constraint system supports separation of noisy and sensitive circuits.

\item \textbf{Prototyping} -- Where rapid iteration and parametric design enable quick exploration of design alternatives.

\item \textbf{Automated Manufacturing} -- Where programmatic generation enables batch production and design automation pipelines.
\end{itemize}

\section{Design Philosophy}
EDADSL follows a declarative-imperative hybrid philosophy:

\begin{itemize}
\item \textbf{Imperative Circuit Definition} -- Components and connections are defined sequentially, allowing natural expression of circuit topology.

\item \textbf{Declarative Physical Constraints} -- Physical requirements are declared as constraints, letting the solver determine optimal placement and routing.

\item \textbf{Separation of Concerns} -- The electrical description (what the circuit does) is separated from the physical description (how it is laid out).

\item \textbf{Constraint Hierarchy} -- A priority system allows designers to express preferences while ensuring critical requirements are met.

\item \textbf{Block Scoping} -- Variables declared inside code blocks (\texttt{if}, \texttt{for}, \texttt{while}, \texttt{do-while}) are local to that block. Functions create their own local scope.
\end{itemize}

